\section{大厂 bets 与创业：为什么研究者下场做公司}

苏煜认为，曾经各家公司在 Agent 上的 bet 差异较大，但现在趋于统一，因为 Anthropic 的路径给行业打了样。Anthropic、OpenAI 都在往 productivity 相关方向收束，Google 拥有强模型和生态位但产品声势不一定匹配，xAI 等也在关注 computer user agent 或 general digital agent。




\subsection{学校、公司与创业的资源结构}

苏煜解释自己为什么从学校转向创业：学校适合 weird ideas、proof of concept 和概念框架；但当 Agent 进入资源密集阶段，需要 GPU、API、强团队和快速工程执行，学校资源结构就不够匹配。创业不是简单为了钱，而是为了继续做研究，只是研究形态变了。

\noindent\begin{tabularx}{\textwidth}{p{0.18\textwidth}p{0.36\textwidth}X}
\toprule
场景 & 适合做什么 & 不适合做什么 \\
\midrule
学校 & 概念框架、低成本验证、长周期 weird ideas & 大规模工程、快速产品迭代、重资源实验。 \\
大厂 & 大模型训练、基础设施、生态整合 & 多方向自由探索可能受组织目标限制。 \\
Startup & 快速试错、聚焦产品、围绕 Agent 建完整系统 & 资源压力大，需要找到明确 wedge 和商业路径。 \\
\bottomrule
\end{tabularx}

\subsection{Conceptual framework 的价值}

访谈最后，苏煜说自己喜欢 build conceptual framework：不是记忆特别好或反应特别快，而是能学很多东西、把它们串起来、看见联系。这也是本期综述最有价值的地方：它把 Agent 的历史、语言、工具、coding、continual learning、world model、产业扩散放进一个统一框架。

\begin{importantbox}{为什么这期值得做成高标准笔记}
许多播客是观点密集但结构松散；本期恰好相反，它本身就在搭概念框架。整理的重点不是复述每个例子，而是保留这套框架：Agent 定义、技术史、Language Agent 栈、边界消融、持续学习瓶颈和产业 bets。
\end{importantbox}

\subsection{本章小结}

Agent 已从 proof of concept 进入资源密集阶段。学校、大厂和 startup 的分工正在变化；能够搭建概念框架、又能调动工程资源的人，会更容易推动下一阶段 Agent。

